\documentclass[12pt,a4paper]{article}
% autocompile publish
\usepackage{math-hse,array}
\title{Лекция 13. Элементы теории вероятностей, часть 2}

\date{26.12.2025}
\begin{document}

В прошлый раз мы обсуждали связь между вероятностью и действием взятия дополнения к событию. Можно аналогичным образом написать формулу для связи вероятности и объединения событий:

\begin{theorem*}
	[Теорема сложения вероятностей]
	Для любых двух событий $A$ и $B$ вероятность того, что хотя бы одно из событий произойдёт равна сумме вероятностей событий $A$ и $B$ минус вероятность их одновременного выполнения:
	$$
		P(A\cup B)=P(A)+P(B)-P(A\cap B).
	$$
\end{theorem*}\section{Условная вероятность}

Рассмотрим ситуацию. Во дворе играют 20 детей: 12 мальчиков и 8 девочек (см.
табл.~\ref{deti}). Дети играют в футбол и в классики. В футбол играют 8 мальчиков и 2 девочки, в классики остальные ребята, т.\,е. ~4 мальчика и 6 девочек. Выберем случайным образом ребёнка среди играющих. С какой вероятностью он играл в футбол? Ответ следует непосредственно из определения вероятности:
$$
	P(F)=\frac{N(F)}{N(\Omega)}=\frac{10}{20}=\frac{1}{2}.
$$

\begin{table}[b]
	\begin{center}
		\begin{tabular}{|c|c|c|c|c|}
			\hline
			& классики & футбол & всего \\
			\hline
			девочки & 6 & 2 & 8 \\
			\hline
			мальчики & 4 & 8 & 12 \\
			\hline
			всего & 10 & 10 & 20 \\
			\hline
		\end{tabular}
		\caption{Дети во дворе}\label{deti}
	\end{center}
\end{table}
Теперь представим себе, что мы видим: выбранный случайным образом ребёнок оказался мальчиком. С какой вероятностью он играл в футбол? Теперь изменилось и количество благоприятных исходов, и общее количество исходов.
$$
	P(F|M)=\frac{N(F\cap M)}{N(M)}=\frac{8}{12}=\frac{2}{3},
$$
где $P(F|M)$~--- вероятность того, что выбранный ребёнок играет в футбол, при условии, что он оказался мальчиком, $n(F\cap M)$~--- количество мальчиков играющих в футбол, $n(M)$~--- общее количество мальчиков. Посмотрим на наше вычисление внимательнее:
$$
	P(F|M)=\frac{\frac{N(F\cap M)}{N(\Omega)}}{\frac{N(M)}{N(\Omega)}}=\frac{P(F\cap M)}{P(M)}.
$$

Здесь $P(F\cap M)$~--- вероятность того, что случайно взятый ребёнок оказался мальчиком, играющим в футбол, $P(M)$~--- вероятность того, что случайно взятый ребёнок оказался мальчиком. Эту формулу можно воспринимать как определение \emph{условной вероятности}: вероятностью события $A$ при условии $B$ называется отношение вероятности их одновременного выполнения к вероятности события $B$.

Важно не путать два понятия: вероятность одновременного выполнения двух событий и вероятность события $A$ при условии $B$. Можно считать, что во втором случае мы рассматриваем не все элементарные исходы, а только те, что благоприятны $B$.

\begin{example}
	При бросании игральной кости события $A$~--- выпадение чётного числа  и $B$~--- выпадение числа, меньшего трёх. Найдём вероятность события $A$ при условии $B$, т.\,е. ~вероятность выпадения чётного числа, при условии, что выпало меньше трех. Мы уже выяснили, что событие $AB$ состоит в выпадении двойки, т.\,е.~ $P(A\cap B)=1/6$,~~ $P(B)=2/6=1/3$, тогда
	$$
		P(A|B) =\frac{P(A\cap B)}{P(B)}=\frac{1/6}{1/3}=\frac{1}{2}.
	$$
\end{example}
\begin{example}
	Случайное испытание: троекратное бросание монетки. Случайное событие $A$~--- выпадение хотя бы двух орлов и $B$~--- выпадение хотя бы одной решки. Найдём вероятность события $A$ при условии $B$, выпадение хотя бы двух орлов при условии выпадение хотя бы одной решки. Мы уже выяснили, что событию $A\cap B$ благоприятствуют три элементарных исхода: «орёл, орёл, решка», «орёл, решка, орёл», «решка, орёл, орёл». Значит, $P(A\cap B)=3/8$, $P(B)=7/8$. Итак:
	$$
		P(A|B) =\frac{P(A\cap B)}{P(B)}=\frac{3/8}{7/8}=\frac{3}{7}.
	$$
\end{example}

\begin{example}
	Рассмотрим ситуацию: вы кидали монетку уже 3 раза, и каждый раз выпадал орёл. С какой вероятностью и на четвёртый раз снова выпадет орёл? Орёл четыре раза подряд выпадает достаточно редко, значит, сейчас скорее выпадет решка, чем орёл, подсказывает нам интуиция. И ошибается! Действительно, вероятность события «выпали четыре решки подряд» составляет $1/16$, т.\,е. достаточно мала. Здесь благоприятным является единственный из 16 возможных исходов: РРРР, ОРРР, РОРР, РРОР, РРРО, ООРР, ОРОР, ОРРО, РООР, РОРО, РРОО, ОООР, ООРО, ОРОО, РООО, ОООО. Условная же вероятность события «в четвёртый раз выпала решка», при условии «первые три раза выпадала решка»  составит $1/2$, т.\,к. ~есть всего два равноправных исхода, в которых первые три раза выпадала решка. Здесь именно выполнение условия является редким событием, а не выпадение решки при условии уже выпавших трёх.
\end{example}

\section{Независимые события}
События $A$ и $B$ называются \emph{независимыми}, если вероятность события $A$ не зависит от того, произошло $B$ или нет. Формально: $P(A|B)=P(A)$.

\begin{example}
	Зависимы ли события «выбранный ребёнок играет в футбол» и «выбранный ребёнок~--- мальчик» из рассмотренного раньше примера?

	Как мы посчитали раньше (см. таблицу~\ref{deti}):

	$P(F)=10/20=1/2$,

	$P(F |M)=8/12=2/3$,

	$P(F)\neq P(F |M)$~--- значит, события не являются независимыми. Это означает, что доля футболистов среди мальчиков больше, чем доля футболистов среди всех детей.
\end{example}

\begin{example}
	 Зависимы ли события: $A$~--- выпадение чётного числа при бросании игральной кости и $B$~--- выпадение числа, меньшего трёх?

	 $P(A)=3/6$,

	 $P(A|B)=1/2$,

	$P(A)= P(A|B)$~--- значит, события независимы. Это означает, что доля чётных чисел среди первых двух такая же, как и среди всех шести возможных.
\end{example}

Мы знаем, что: $P(A|B) =P(A\cap B)/P(B)$. Значит, условие независимости событий можно записать так: $P(A)=P(A|B)=P(A\cap B)/P(B)$. Откуда получаем:

$$P(A\cap B)=P(A)\times P(B)$$

Это важное свойство независимых событий, часто его берут за основное определение.

Иногда вопрос о зависимости двух событий можно рассматривать как математическую задачку, как мы делали выше. Бывают же ситуации когда, наоборот, из природы событий ясно, что они независимы, тогда их независимостью можно пользоваться для нахождения вероятности одновременного выполнения событий.

\begin{example}
	Монетку бросили три раза. С какой вероятностью выпало три решки? Эту задачу мы уже решили, подсчитав общее количество элементарных исходов и количество благоприятных исходов. Посмотрим теперь на этот вопрос так: интересующее нас событие $A$ состоит в одновременном выполнении трёх событий: $B$~--- в первый раз выпала решка, $C$~--- во второй раз выпала решка, $D$~--- в третий раз выпала решка. Эти события независимы (если, конечно, мы не заинтересованы в результате и бросаем монетку честно). Значит,
	$$
		P(A)=P(B\cap C\cap D)=P(B)\times P(C)\times P(D)=\frac{1}{2}\times\frac{1}{2}\times\frac{1}{2}=\frac{1}{8}.
	$$
\end{example}

\begin{example}
	 Монетку бросили три раза. Найти вероятность события $A$~--- выпал хотя бы один орёл. На первый взгляд, мы не можем применять в этой задаче соображения независимости событий, ведь «выпал хотя бы один орёл»~--- это значит, что мог выпасть именно один орёл, в любой из разов, могло выпасть два орла, могло выпасть три орла, т.\,е.~ само событие устроено сложно. В чем состоит дополнительное событие $B=\Omega \verb|\| A $? Другими словами, в каких случаях событие A не выполняется? Когда не выпало ни одного орла, т.е. выпали три решки. Вероятность этого события мы только что нашли: $P(B)=1/2\times1/2\times1/2=1/8$. Значит, $P(A)=1-P(B)=7/8$.
\end{example}
\end{document}
